On réalise une expérience en introduisant, à l'instant $t_{0}=0$, une masse de
zinc en poudre de valeur ${m(\ce{Zn})=\qty{1.0}{\g}}$ dans un ballon contenant
le volume $V=\qty{40}{\mL}$ d'une solution aqueuse $\mathrm{(S)}$ d'acide
chlorhydrique ($\ce{H3O^{+}_{(aq)} + Cl^{-}_{(aq)}}$) de concentration molaire
$C_{A}=\qty{0.5}{\mol\per\L}$. Les ions $\ce{H3O^{+}_{(aq)}}$ réagissent avec le
zinc $\ce{Zn_{(s)}}$ suivant la réaction chimique d'équation~:
\[
    \ce{2H3O^+_{(aq)} + Zn_{(s)} -> H2_{(g)} + Zn^2+_{(aq)} + 2H2O_{(\ell)}}.
\]

La mesure du volume de dihydrogène formé permet le suivi de l'évolution
temporelle de l'avancement $x$ de la réaction et de tracer le graphe $x=f(t)$.

\begin{donnees}
    \begin{itemize}
        \item $M(\ce{Zn})=\qty{65.4}{\g\per\mol}$
    \end{itemize}
\end{donnees}

\begin{figure}[H]
    \centering
    \includegraphics[width=0.5\textwidth]{2025-11-27_110541-}
\end{figure}

\begin{enumerate}
    \item Calculer les quantités de matière $n_{0}(Zn)$ et $n_{0}(\ce{H3O+})$,
          présentes initialement dans le mélange réactionnel.
    \item Recopier, sur votre copie, le tableau d'avancement de la réaction
          chimique et le compléter.

          \begin{table}[H]
              \flushright
              \setlength{\arrayrulewidth}{0.8pt}
              \renewcommand{\arraystretch}{1.5}
              \begin{tabularx}{0.95\textwidth}{|p{3.3cm}|
                    >{\centering\arraybackslash}p{3.3cm}|C|C|C|C|C|}
                  \hline
                  \multicolumn{2}{|c|}{Équation chimique} &
                  \multicolumn{5}{c|}{%
                      \ce{2H3O^{+}_{(aq)}} \hskip2mm $+$ \hskip3mm
                      \ce{Zn_{(s)}} \hskip3mm
                      $\longrightarrow$ \hskip2mm \ce{H2_{(g)}} \hskip3mm
                      $+$ \hskip3mm
                      \ce{Zn^{2+}_{(aq)}} \hskip3mm $+$ \hskip3mm
                      \ce{2H2O$_{(\ell)}$}} \\
                  \hline
                  État du système & Avancement (\unit{\mol}) &
                  \multicolumn{5}{c|}{Quantités de matière (\unit{\mol})} \\
                  \hline
                  État initial & $x=0$ &  &  &  &  & excès \\
                  \hline
                  État intermédiaire & $x$ &  &  &  &  & excès \\
                  \hline
                  État final & $x_{f}$ &  &  &  &  & excès \\
                  \hline
              \end{tabularx}
          \end{table}

    \item Identifier le réactif limitant. Justifier.
    \item Déterminer graphiquement~:
          \begin{enumerate}
              \item la valeur du temps de demi-réaction $t_{1/2}$.
              \item la valeur de la vitesse volumique de réaction, en unité
                    (\unit{\mol\per\L\per\s}), à l'instant $t=\qty{400}{\s}$,
                    sachant que le volume du mélange réactionnel est
                    $V=\qty{40}{\mL}$.
          \end{enumerate}
    \item Interpréter qualitativement la variation de la vitesse volumique de
          cette réaction.
    \item Pour accélérer la réaction précédente, on recommence l'expérience en
          utilisant la même masse de zinc $m(\ce{Zn})=\qty{1.0}{\g}$ et le
          volume $V=\qty{40}{\mL}$
          d'une solution aqueuse ($S^{\prime}$) d'acide chlorhydrique de
          concentration molaire $C_{A}{ }^{\prime}=\qty{1}{\mol\per\L}$.
          \begin{enumerate}
              \item Citer le facteur cinétique qui est à l'origine de
                    l'accélération de la réaction.
              \item Le temps de demi-réaction $t_{1/2}$ va-t-il augmenter ou
                    diminuer~? Justifier.
          \end{enumerate}
\end{enumerate}


